\documentclass[../main.tex]{subfiles}

\begin{document}



\textbf{II) \underline{Using a gluing theorem}}
\newline

The goal of this part is to recollect a sequence of regular pages from two rigid pages having a common asymptotic orbit, being negative for one and positive for the other. Take $2$ such curves $u_+$ and $u_-$ respectively. \\

We start by computing the Fredholm indices of such curves. We compute their 1st Chern number and Conley-Zehnder index in convenient trivializations, namely the one defined before. Near a radial orbit choose the natural trivialization induced by the coordinates $(\rho,\phi,\theta)$. In these coordinates, we assumed that the Conley-Zehnder index of the orbits $\gamma_i$ is $+1$ for a positive orbit, or $-1$ for a negative one.\\

Near a broken component, we also choose the trivialization given by the coordinates defined before. Now in this trivialization the Conley-Zehnder index is $0$, because the orbit is positive hyperbolic, and the eigenvectors have to turn by an angle less than $\frac{\pi}{4}$. This is because if such a vector were to turn more for the linearized Reeb flow, this would induce a transverse Reeb flow on a rigid page with the wrong orientation (recall the rigid pages have equation $\{\phi = \text{Cst}\}$, so they do not wind around the Reeb orbit in this trivialization). Hence by definition of the Conley-Zehnder index on positive hyperbolic orbits, it has to be $0$.\\
    
Finally, take a lifted rigid page $u$ : we claim its relative normal 1st Chern number (see \cite{wendl_automatic_2009} for definitions and details), denoted $c_i^{\tau}(N_u)$, is $0$. Indeed, away from binding components one can take $R_\alpha$ as a positive normal vector field, and near the binding components one can interpolate (by positive transverse vector fiedls) $R_\alpha$ with $\partial_\phi$ or $-\partial_\phi$ depending on the positivity of the page. $\pm\partial_\phi$ is a positive transverse vector field and is constant "at infinity" in the trivialization considered. Moreoever up to an homotopy of the plane field $\xi$, we can assume that $\partial_\phi\in\xi$, so we have a global section of a generalized normal bundle of $u$ (see \cite{wendl_automatic_2009}). Hence we get $c_1^\tau(N_u) = 0$, with $\tau$ the trivialization.  \\

Hence we can compute the Fredholm index of our curves : denote by $k_+$ the number of positive ends, $k_-$ the number of negative ends. Then the index is :

$$\text{Ind}(u) = \chi(\dot \Sigma) + c_1^\tau(N_u) + (k_+ + k_- - 1) = (2-k_+ - k_-) + (k_+ + k_- -1) = 1$$

Hence the index is always $1$ for this kind of curves. Note that positive and negative orbits play the same role in the computation of the index, because negative orbits have opposite Conley-Zehnder index than positive orbits.\\

\begin{remark}
   As noticed at the end of the last part, we can apply exactly the same construction for regular pages, considering only asymptotic radial binding components. In that case, the very same index computation raises : $\text{Ind}(u) = 2$. Hence our construction allows us in fact to lift an arbitrary finite number of curves of the broken book, rigid or regular, as a $1$-parameter family of holomorphic curves in the symplectization with Fredholm index $1$ or $2$. 
\end{remark}

Now before going through the gluing process, we use \cite{wendl_automatic_2009} to show that this index computation actually gives the Fredholm regularity of all those curves, thanks to some "automatic transversality" phenomenon, which is very specific to dimension $4$. Applying Proposition 2.2 of this article in our current setup shows that if $D$, the linearized Cauchy-Riemann operator, verifies $\text{Ind(D)}> 0$ then it is surjective. This is indeed the case here, hence all our lifted curves are Fredholm regular, and applying a generic perturbation of $J$ does not destroy the curves (up to isotopy relative to the asymptotic orbits). This allows us to suppose that our $J\in\J(\alpha)$ is generic inside some open set of $\J(\alpha)$. \\

With this in hand, we can finally the gluing theorem from \cite{hutchings_gluing_2007}, which tells us that under genericity assumptions that are verified here, we can glue $u_+$ and $u_-$ to get a sequence of index $2$ curves $u_\delta$ which are $\delta$-close from breaking, and which breaks on the $1$-level holomorphic building composed of $u_-$ and $u_+$ glued together. 

\begin{remark}
    The theorem actually gives a sequence of curves converging to a building with upper level $u_+$ and lower level $u_-$, and maybe some ramified covers of trivial cylinders between them. However here $u_+$ and $u_-$ are somewhere injective curves asymptotic to $\gamma_b$ which is a simply covered Reeb orbit, hence no such cylinder can appear. In particular, for any $n$, $[u_n] = [u_-] + [u_+]$ in $H_2(M,\cup_i\gamma_i)$.
\end{remark}

\begin{remark}
    In fact, the computation of the quantity $\#G(u_-,u_+)$ which is actually given by the theorem raises the value $1$, which means that there are $1$ mod $2$ ways to glue our rigid curves. We will show in the next part that there is indeed one, corresponding to regular pages.
\end{remark}

                
Finally, notice that a little bit of ECH theory gives us embeddedness of the curves obtained in the previous sequence. Indeed, $u_-$ and $u_+$ are both embedded  hence $I(u_-) = I(u_+) = \text{Ind}(u_-) = \text{Ind}(u_+) = 1$, where $I(u)$ is the ECH index. Now by additivity property of this index, if $u$ is some curve "close to breaking", with $[u] = [u_-] + [u_+]$, then $I(u) = 2 = \text{Ind}(u)$, hence $u$ is embedded. We refer to \cite{hutchings_lecture_2014} for details on the ECH index.

\begin{remark}
    In the last part we will make use of Siefring's intersection theory rather than ECH, and the embeddedness property of the curves could also be proved quickly with the adjonction formula in that case.
\end{remark}
 

\textbf{III) \underline{Recollecting the full foliation}}
\newline

In this last part we finally lift the full foliation. All the curves will be denoted with the notation $u = (a,f)$, with $a$ the symplectization coordinate and $f$ the projection on $M$.\\

First, consider $(M,\alpha)$ as before, with supporting BBD verifying the conditions of part $1$. Then $M$ is divided in several "sectors" delimited by $4$ rigid pages. To define such a sector, take any regular page and notice than pushing it by a reparameterization of the Reeb flow allows to recollect the nearby pages, as long as they are transverse to the Reeb flow. If at some point the foliation cannot be obtained that way, it means that some part of the foliation is not transverse to the Reeb flow anymore, which can only happen at a broken binding component. Hence if $F$ is a regular page, we have $2$ possitilities : either there is no problem of transversality whatsoever, and in this case $F$ is a Birkhoff section, and we have a nice open book decomposition ; or we get a sector of $M$ on which the foliation is $]0,1[\times M$, and which has boundary components consisting of $2$ rigid pages, asymptotic to the same broken binding component with ends of opposite sign. Note that because of the local model around a broken binding component, the boundaries at $\{0\}$ and the boundaries at $\{1\}$ cannot be the same. Finally, because any point in $M$ belongs to a page, we get that $M$ can be divided into a finite number of such sectors, "glued" to each other along rigid pages and binding components.\\

We consider in the following that there is at least one broken component in the BBD, hence at least one rigid page . Thus consider any such sector $S$. We denote as $F$ some regular page, and $F_-^0, F_+^0,F_-^1,F_+^1$, the four rigid pages delimiting the sector. These four pages can be lifted as  $u_-^0, u_+^0,u_-^1,u_+^1$
$J$-holomorphic curves for some $J\in \J(\alpha)$. Denote $A = [F] = [u_-] + [u_+] \in H_2(M,\cup_i\gamma_i)$ where the $\gamma_i$ are the radial binding components, which are also the asymptotic orbits of $F$. \\
    
Take $u_-^0$ and $u_+^0$ with broken orbit $\gamma_b^0$ ; the previous part gives us $(u_n)$ index-2 curves such that $[u_n] = A$ and $u_n\rightarrow u_\infty = (u_-^0,u_+^0)$ in the classical SFT-comptactness sense. We know that $u_n$ are embedded curves in the symplectization, let us show that their image belong to $\R\times S$. To do that, we will make use of the intersection product defined by Siefring \cite{siefring_intersection_2011}, we refer to \cite{wendl_contact_2019}, Appendix C, for the main results we will use. The main point is that there exist a pairing $u*v$ for $u,v$ J-holomorphic curves, which only depends of the relative homology class of $u$ and $v$, and does not depend on any choice of trivialization. \\

Take $u = u_n$ for some $n$ (that will be needed large enough in what follows) and let us compute $u*u$. The adjunction formula for this pairing is (see \cite{wendl_contact_2019} for the different notations) :

$$u*u = 2[\delta(u) + \delta_\infty] + c_N(u) + [\bar{\sigma}- \#\Gamma_u]$$

Here, because all the asymptotic orbits are simply covered, we have $[\bar{\sigma}- \#\Gamma_u] = 0$ and $\delta_\infty = 0$. $\delta(u)$ is the algebraic count of self-intersections of $u$, which is also $0$. Finally, by the index computations already made before, the normal Chern number $c_N(u)$ is also $0$. Hence $u*u = 0$. A similar computation raises the same result for the lifted rigid curves. \\

Now, take any rigid curve $v$, and compute : $u*v = u\cdot v + \iota_\infty(u,v)$, where the first term is the algebraic count of actual intersections, and $\iota_\infty(u,v)$ is the "hidden intersections" that appear after a small perturbation of $v$ given by some trivialization. However, we know that there is no such hidden intersections at radial binding components by definition of the local model (we can perturb $v$ close to the Reeb orbit following $\partial_\phi$, and counts the intersections that appear, which is $0$). Hence $u*v = u\cdot v$ is the actual number of intersections (which is the absolute number, by positivity of intersections). Now take $u_0$ a lift of a regular page in $S$ far from the rigid pages. We have that $u_0*v = 0$ because $u_0$ does not intersect any lift of rigid pages. This modifies the almost complex structure locally, we up to chosing $u_0$ far enough from the rigid pages and $n$ large enough, it does not affect $u_n$. Hence $u*v = [u]*v = u_0*v = 0$. \\

Now suppose $u$ is not in $\R\times S$. Then if we suppose $n$ is large enough, we can look at $u$ in the local hyperbolic model around $\gamma_b$. Denote as $\U_i$ the considered neighborhood. If $\text{Im}(u)\cap (\R\times \U_i) $ is not entirely in $\R\times S$, then by definition of the local model, $u$ has to intersect a lift of a rigid page, which contradicts the previous intersection considerations. Away from the ends, one can parameterize a neighborhood of $f_-^0$ with the flow of $\pm R_\alpha$. Then $f$ is contained in $\delta > t > 0$ for some $\delta > 0$ small, because $u$ cannot intersect $u_-^0$ or its translations. Hence $f$ has its image inside $S$.\\

Moreover, we also know that the projection of $u$ on $M$ is an embedding ($u$ is "nicely embedded" with the terminology of \cite{wendl_automatic_2009}) because $u_-^0$ and $u_+^0$ are nicely embedded, and by definition of the SFT-convergence. More precisely, self-intersections are not possible for $n$ large enough, because of the $\mathcal{C}^0$ convergence of the projections, and one can also shows that $f$ is an immersion. Let us precise a little bit this point, as this notion will be useful to study the compactness of the moduli space.\\

The notion of nicely embedded curves was introduced by Wendl in \cite{wendl_automatic_2009} in order to study curves in finite energy foliations with partcularly nice properties implying a good behavior for limits of sequences of such curves. In the case of a curve $u$ in a symplectization, being nicely embedded is equivalent to either being a trivial cylinder, or to verify the following conditions (see \cite{cioba_nicely_nodate}) :

\begin{enumerate}
    \item $u$ is somewhere injective.
    \item $u*u\leq 0$.
    \item $\delta(u)+\delta_\infty(u) = 0$.
\end{enumerate}

From there, it is fairly straightforward to deduce that given a nicely embedded curve $u$, all the curves in the same connected component as $u$ in the moduli space of all curves are nicely embedded too. Moreoever, it is always true in that case that they are all either identical or disjoint, as well as Fredholm regular, which we already proved before in our particular case.\\

Now denote by $\mathcal{M}(A,J)$ the connected component of $u$ in the moduli space of asymptotically cylindrical curves with homology class $A$ modulo translation in the symplectization. This is a one-dimensional manifold by computation of the Fredholm index and regularity of the curves. Let us show that its boundary can be identified with the union of $2$-level buildings : $(u_-^0,u_+^0)\cup(u_-^1,u_+^1)$.\\ 

The proof of that relies the convergence properties of nicely embedded curves in symplectizations, as studied in \cite{wendl_compactness_2008}. First note that $\mathcal{M}(A,J)$ has non-empty boundary. If not, we could look at the evaluation map : $$ev : \mathcal{M}_1(A,J)\rightarrow M$$ where $\mathcal{M}_1(A,J)$ is the moduli space of curves in $\mathcal{M}(A,J)$ with one marked point, modulo reparameterization. Notice that this space has compactification that we can write : $\overline{\mathcal{M}_1(A,J)} = \mathcal{M}_1(A,J)\cup (\bigcup\gamma_i)$ where we identified the asymptotic orbits of $u$ with the trivial cylinders covering them. The evaluation map extends trivially on this space, and is smooth with degree one, hence by every point of $M$, we can find a curve in the homology class of $u$ which passes through that point. This is a contradiction because $u$ cannot cross any rigid page by intersection property. Hence $\mathcal{M}(A,J)$ is a connected $1$-manifold with boundary $u_\infty^0\cup u^1_\infty$. By the same intersection argument, any regular curve of $\mathcal{M}(A,J)$ has image inside $\R\times S$.\\

Suppose $v_n$ is a sequence of $\mathcal{M}(A,J)$ converging to $u_\infty$ (meaning $u_\infty^0$ or $u_\infty^1$). By genericity, we can assume that any somewhere injective curve is Fredholm regular, hence has positive Fredholm index if it is not a trivial cylinder. We know that $v_n$ is a sequence of stable, nicely embedded curves because they are in the same moduli space connected component as $u$. We will now use the main theorem from \cite{wendl_compactness_2008} to conclude that $(v_n)$ can only degenerate in very convenient ways. We recall here the main features of this theorem that we will need here :

\begin{definition}{(Definition 1.5 of \cite{wendl_compactness_2008})\\}
    A stable J-holomorphic building in $\R\times M$ is said to be nicely embedded if :
    \begin{enumerate}
        \item It has no nodes.
        \item Any connected component is a trivial cylinder or is nicely embedded.
        \item The projections on $M$ of $2$ disctincts components are identical or disjoint.
        \item Any breaking orbit is even and simply covered or even and doubly covered and bad (see \cite{} for the definition of a bad orbit).
    \end{enumerate}
\end{definition}

\noindent The following result works only in the setting of a symplectization and not in an arbitrary symplectic cobordism.

\begin{theoreme}{(SFT Convergence for nicely embedded curves, Theorem 1 of \cite{wendl_compactness_2008})\\}
    Suppose that $(u_n)$ is a sequence of nicely embedded curves with $u_n : \dot\Sigma \rightarrow \R\times M$ and suppose that $u_n\rightarrow u$ where $u$ is a stable $J$-holomorphic building. Then $u$ is nicely embedded.
\end{theoreme}

Applying this theorem, we see that no bubbling or nodal degeneration can happen in our case. \\

Now suppose $v_n$ breaks into a $N$-level holomorphic building, which is therefore nicely embedded. Then we claim that the only orbit that can be involved in the intermediate levels is $\gamma_b^0$ or $\gamma_b^1$. Indeed, $v_n$ breaks on orbits in $\overline{S}$, which are all binding components. The previous theorem tells us that the sequence can only break on even Reeb orbits or double covers of odd hyperbolic orbits, however there are no odd hyperbolic orbits in the binding so it can only break on even Reeb orbits. Then because radial binding components are supposed elliptic and odd, we get that the breaking can only happen on $\gamma_b^0$ or $\gamma_b^1$. \\

If we now assume that $(v_n)$ breaks on $\gamma_b^0$, then it cannot break on $\gamma_b^1$ as this would force $u_\infty$ to cross $u$ at some point which would contradict the self-intersection computation (as explained in \cite{wendl_contact_2019}, the intersection product extend to holomorphic buildings, and by proposition $C.6$, $0 = u*u = v_n*u\rightarrow u_\infty*u= 0$, which implies that no level intersects $u$). Hence $(v_n)$ can only break on $\gamma_b^0$. We must show that the limit building is then $(u_-^0,u_+^0)$. First look at the local model around $\gamma_0^b$. If a non-constant component of the limit building were distinct from $u_0^+$ or $u_-^0$, because it has to be in $S$, that means it crosses some curve $u_n$ for $n$ large enough. By the properties of the intersection product for holomorphic buildings, (\cite{wendl_contact_2019}) this is not possible. Indeed, the limit building has intersection $0$ with $u_n$ by continuity property of the intersection product, and by superadditivity this implies that the intersection product of $u_n$ with all non-trivial components of the building is $0$. Hence by continuation principle, the limit building is $(u_-^0, ..., u_+^0)$ where "..." stands for an unknown number of ramified covers of trivial cylinders over $\gamma_b^0$. However, by the previous theorem, the building has to be nicely embedded, hence has no multiply covered components (we could also use the same argument as during the gluing process, i.e these covers cannot exist in a such a building anyway as $\gamma_b^0$ is simply covered by $u_-^0$ and $u_+^0$). Note that the exact same reasonning holds if $(v_n)$ breaks on $\gamma_b^1$. \\

Therefore we know that if the sequence $(v_n)$ breaks, it has to be on $\gamma_b^0$ or $\gamma_b^1$ and the limit building has to be $(u_-^0,u_+^0)$ or $(u_-^1,u_+^1)$. Moreover it is easy to see that the $2$ components of $\partial\overline{\mathcal{M}(A,J)}$ cannot be the same (S is divided in $2$ parts by $f$, and each boundary component is in one of these parts so they cannot be the same) from which we deduce that $\partial\overline{\mathcal{M}(A,J)} = (u_-^0,u_+^0)\cup(u_-^1,u_+^1)$. \\

Looking at the moduli space of marked curves again, we get an evaluation map : $ev :\overline{\mathcal{M}_1(A,J)}\rightarrow \overline{S}$ which has degree one (here the degree is defined between compact manifolds with boundary). Hence it is surjective, and any point in $S$ belong in the image of some curve (more precisely to the projection of a 1-parameter family of curves in the symplectization) of $\mathcal{M}(A,J)$. Moreover, 2 such curves are either disjoint or the same, hence exactly one curve passes through every point, and $S\simeq \text{Im}(f)\times ]0,1[$ (and even more : $\overline{S} \simeq \text{Im}(f)\times [0,1]$ topologically, by considering $u_-^0\#u_+^0$ and $u_-^1\#u_+^1$). \\

Now we can apply this construction to all sectors in the manifold. We thus get a finite energy foliation of $M$ as wanted. The last thing that remains to prove is that this foliation is isotopic to the initial broken book decomposition. This is fairly straightforward as we know the Reeb vector field is transverse to both foliations (see \cite{cioba_nicely_nodate} Lemma $2.1.15$ for a proof that nicely embedded curves have a projection transverse to the Reeb flow). Chose some sector $S$, and take some Reeb orbit $\gamma$ which is not an asymptotic orbit of any page in the sector (rigid or regular), with $\text{Im}(\gamma)\subset \overline{S}$, and such that $\gamma(0)$ and $\gamma(1)$ belong to different rigid pages of $\overline{S}$. Denote as $\F$ the foliation of the broken book decomposition, and $\F'$ the foliation induced by the projection of the holomorphic curves. Then define : 
$$\pi : S \rightarrow [0,1]$$

\noindent by associating to a point $x$ the time $t$ for which $\gamma(t) \in \F_x$ with $\F_x$ the page containing $x$. Because the Reeb vector field is transverse to $\F$, this is well defined and smooth. Consider $h : S \rightarrow \R$ such that $d_x\pi(R_\alpha) = h\partial_t$ with $\partial_t$ corresponding to the $[0,1]$-coordinate. In the same way one can define $\pi'$ and $h'$ associated to $\F '$. Now consider the map :  

\begin{align*}
    F : \quad  &S\times [0,1]\longrightarrow S\\
    &(x,s) \longmapsto\phi_{\pi(x)}^{\frac{s}{h}R_\alpha  }\circ\phi^{\frac{1-s}{h'}R_\alpha}_{\pi'(x)}\circ\phi_{\pi(x)}^{-\frac{R_\alpha}{h}}(x)
\end{align*}

This gives the isotopy we wanted. Because $F$ is the identity on the rigid pages and on the binding orbits, we can compose all the isotopies of the different sectors to get a $\mathcal{C}^0$-isotopy between the projection of the finite energy foliation and the broken book decomposition. This isotopy can be chosen smooth away from the binding orbits up to reparametrizing the previous path $\gamma$, in such a way that $\pi$ and $\pi'$ are "flat" on the rigid pages. 


\begin{remark}
    Note that our construction also allows to lift regular open book decompositions as a finite energy foliation, hence giving a new slightly different proof of this result in the particular case of elliptic binding orbits of Conley-Zehnder index 1 (which also allows negative binding components). This corresponds to the case where there is no rigid page in the decomposition. In order to adapt the proof, first lift a page $F$ as a holomorphic curve $u$, and look at the associated moduli space as before. It turns out that this time the curve cannot break, because as shown before it cannot break on a binding component, and all the other closed orbits are crossing $F$ transversally, so no sequence of curves in the same connected component of the moduli space as $u$ can break on such an orbit because that would imply the existence of intersections between $u$ and curves close to breaking. So the moduli space is compact with no boundary, hence a circle, and the curves foliate the whole space by similar arguments than above, which concludes.
\end{remark}




\end{document}